\input{../header}

\begin{document}
%


%\onehalfspacing
\allowdisplaybreaks
%##################################################################
\section{PS\#13 - Taylor series - \red{Answer key} }

\everymath{\displaystyle}
\begin{enumerate}[leftmargin=0pt]


\item (\href{https://faculty.gvsu.edu/boelkinm/Home/ACS/sec-8a-4-taylor-series.html#act-8a-4-1}{AC Activity 8a.4.2}) Consider the infinite Taylor series given by $T(x) = \sum_{k = 1}^{\infty} \frac{1}{k \cdot 2^k}(x-1)^k.$
\begin{enumerate}
    \item As described in the statment of the Ratio Test, let $r_n(x)$ be the ratio of the $(n+1)$st term of $T(x)$ to the $n$th term of $T(x)$. Find the simplest formula that you can for $r_n(x)$.

    \begin{red}
        $r_n(x) = \dfrac{n}{n+1} \cdot \dfrac12 \cdot (x-1).$
    \end{red}
    \item Let $r(x) = \displaystyle\lim_{n\to\infty} r_n(x)$. Evaluate this limit to find the simplest formula you can for $r(x)$.

    \begin{red}
        $r(x) = \dfrac12 (x-1).$
    \end{red}
    \item For what values of $x$ is $|r(x)| < 1$? \begin{red} $-1 < x < 3$.
    \end{red}
\end{enumerate}

\item In this problem we're going to explore the Taylor series for $f(x) = \sqrt{x}$ centered at $a = 1$.
\begin{enumerate}
    \item Find the Taylor polynomials for $f(x) = \sqrt{x}$ up through order 5, centered at $a=1$.

    \begin{red}
        $T_5(x) = 1 
        + \dfrac12 (x-1) 
        - \dfrac{1}{2!\cdot 2^2}(x-1)^2
        + \dfrac{1\cdot 3}{3!\cdot 2^3}(x-1)^3 
        - \dfrac{1\cdot3 \cdot 5}{4!\cdot 2^4}(x-1)^4
        + \dfrac{1\cdot3 \cdot 5\cdot 7}{5!\cdot 2^5}(x-1)^5$.
    \end{red}
    \item Spot the pattern in the coefficients.

    \begin{red}
        $c_k = (-1)^{k+1}\dfrac{1\cdot 3 \cdot \ldots \cdot (2k-3)}{k!\cdot 2^k}.$
    \end{red}
    \item Use the pattern you found to do the ratio test.
    \begin{red}
    \begin{align*}
        |r_n(x)|
        &= \dfrac{(2n-1)}{(n+1)\cdot 2}\cdot |x-1| \\
        |r(x)| &= |x-1| \lim_{n\to \infty} \dfrac{2n-1}{2n+2} = |x-1|\cdot 1
    \end{align*}
    \end{red}
    
    \item What is the ``interval of convergence'' -- the $x$-values for which the Taylor series will converge? Why does your answer make sense based on the graph of the original function $f(x) = \sqrt{x}$?

    \begin{red}
        This series will converge whenever $|r(x)| = |x-1| < 1$ -- that is, for $0 < x < 2$. This makes sense because the graph of $\sqrt{x}$ stops at $x=0$.
    \end{red}
    \item Use $T_5(x)$ to find a good fraction approximation for $\sqrt{2} \approx 1.414\ldots$

    \begin{red}
        \begin{align*}
            T_5(2) &= 1 
        + \dfrac12 (2-1) 
        - \dfrac{1}{2!\cdot 2^2}(2-1)^2
        + \dfrac{1\cdot 3}{3!\cdot 2^3}(2-1)^3 
        - \dfrac{1\cdot3 \cdot 5}{4!\cdot 2^4}(2-1)^4
        + \dfrac{1\cdot3 \cdot 5\cdot 7}{5!\cdot 2^5}(2-1)^5 \\
        &= 1 
        + \dfrac12 
        - \dfrac{1}{2!\cdot 2^2}
        + \dfrac{1\cdot 3}{3!\cdot 2^3}
        - \dfrac{1\cdot3 \cdot 5}{4!\cdot 2^4}
        + \dfrac{1\cdot3 \cdot 5\cdot 7}{5!\cdot 2^5} \\
        &= \dfrac{365}{256} = 1.42578125.
        \end{align*}
        Not bad!
    \end{red}
\end{enumerate}

    

\end{enumerate}
\end{document}